% GENERATED FILE. Edit canonical lessons, then run npm run generate:books.
% canonical-source-hash: fnv1a64:7ca87b51934edb5c
% canonical-lessons: HI-C281-janachahiye, HI-C281-atah, HI-W281-sign-visarga, HI-C281-isliye, HI-R281-advice-and-reasons

\chapter{Should, Must, and So}
\label{ch:hi-should-must-and-so}
\hlchaptermodality{281}

\begin{chapteropening}
\textbf{By the end of this chapter:} \emph{I can say what someone should do (jānā cāhiye) and has to do (jānā paṛegā), join a cause to its result with isliye and to, and read and write \hi{ः} in ataḥ.}

Complete HI-R281-advice-and-reasons: give advice, say what must be done, and draw the result.
\end{chapteropening}

\section[\texorpdfstring{\hi{तुम्हें} \hi{जाना} \hi{चाहिए} (tumhẽ jānā cāhiye)}{तुम्हें जाना चाहिए (tumhẽ jānā cāhiye)}]{\hi{तुम्हें} \hi{जाना} \hi{चाहिए} (tumhẽ jānā cāhiye) — should, and have to}

\label{lesson:HI-C281-janachahiye}



\begin{quote}
Say \textquotedblleft{}I need tea\textquotedblright{}. (\textbf{\hi{मुझे} \hi{चाय} \hi{चाहिए}}.) Then \textquotedblleft{}I have to go\textquotedblright{}, the way you take leave. (\textbf{\hi{जाना} \hi{है}}.)
\end{quote}

\subsection*{You'll want to know: should}
\begin{quote}
\textbf{\hi{तुम्हें} \hi{आराम} \hi{करना} \hi{चाहिए।}}
\end{quote}

\begin{quote}
\emph{tumhẽ ārām karnā cāhiye.} — \textquotedblleft{}You should rest.\textquotedblright{}
\end{quote}

\textbf{\hi{चाहिए}} after a thing is \textquotedblleft{}is needed\textquotedblright{}. After a verb in \textbf{-\hi{ना}} it is \textquotedblleft{}should\textquotedblright{}. The one who should takes \textbf{\hi{को}} (\textbf{\hi{आपको}}), or a pronoun's short form of it (\textbf{\hi{मुझे}}, \textbf{\hi{तुम्हें}}, \textbf{\hi{उसे}}), and \textbf{\hi{चाहिए}} does not change.

\begin{itemize}
\raggedright
  \item \textbf{\hi{आपको} \hi{डाक्टर} \hi{के} \hi{पास} \hi{जाना} \hi{चाहिए।}} (\emph{āpko ḍākṭar ke pās jānā cāhiye.}) — You should go to the doctor.
  \item \textbf{\hi{हमें} \hi{पानी} \hi{पीना} \hi{चाहिए।}} (\emph{hamẽ pānī pīnā cāhiye.}) — We should drink water.
\end{itemize}

\textquotedblleft{}Should not\textquotedblright{} puts \textbf{\hi{नहीं}} before the verb: \textbf{\hi{तुम्हें} \hi{ऐसा} \hi{नहीं} \hi{करना} \hi{चाहिए}} (\emph{tumhẽ aisā nahīṃ karnā cāhiye}), \textquotedblleft{}you shouldn't do that\textquotedblright{}.

\subsection*{You'll want to know: have to}
\begin{quote}
\textbf{\hi{कल} \hi{मुझे} \hi{दफ़्तर} \hi{जाना} \hi{पड़ेगा।}}
\end{quote}

\begin{quote}
\emph{kal mujhe daftar jānā paṛegā.} — \textquotedblleft{}Tomorrow I will have to go to the office.\textquotedblright{}
\end{quote}

\textbf{\hi{पड़ना}} (\emph{paṛnā}) after a verb in \textbf{-\hi{ना}} says there is no choice: something outside you makes it so. The one it falls on takes \textbf{\hi{मुझे}} again.

\begin{itemize}
\raggedright
  \item \textbf{\hi{मुझे} \hi{रोज़} \hi{जाना} \hi{पड़ता} \hi{है।}} (\emph{mujhe roz jānā paṛtā hai.}) — I have to go every day.
  \item \textbf{\hi{मुझे} \hi{जाना} \hi{पड़ा।}} (\emph{mujhe jānā paṛā.}) — I had to go.
  \item \textbf{\hi{मुझे} \hi{जाना} \hi{पड़ेगा।}} (\emph{mujhe jānā paṛegā.}) — I will have to go.
\end{itemize}

Three shades, then: \textbf{\hi{जाना} \hi{है}}, which you say as you leave, is a plan; \textbf{\hi{जाना} \hi{चाहिए}} is advice; \textbf{\hi{जाना} \hi{पड़ेगा}} is a necessity.

\subsection*{Your turn}
Say these aloud:

\begin{itemize}
\raggedright
  \item \textquotedblleft{}You should go to the doctor\textquotedblright{} — \emph{āpko ḍākṭar ke pās jānā cāhiye}
  \item \textquotedblleft{}You shouldn't do that\textquotedblright{} — \emph{tumhẽ aisā nahīṃ karnā cāhiye}
  \item \textquotedblleft{}Tomorrow I will have to work\textquotedblright{} — \emph{kal mujhe kām karnā paṛegā}
\end{itemize}

\begin{tcolorbox}[breakable,colback=teal!4,colframe=teal!35!black,title={Before you move on}]
\textquotedblleft{}You should rest\textquotedblright{}? (\textbf{\hi{तुम्हें} \hi{आराम} \hi{करना} \hi{चाहिए}}.) Does \textbf{\hi{चाहिए}} change? (No.) \textquotedblleft{}I had to go\textquotedblright{}? (\textbf{\hi{मुझे} \hi{जाना} \hi{पड़ा}}.) \textquotedblleft{}I will have to go\textquotedblright{}? (\textbf{\hi{मुझे} \hi{जाना} \hi{पड़ेगा}}.) Which one says there is no choice? (\textbf{\hi{पड़ना}}.)
\end{tcolorbox}
\section[\texorpdfstring{\hi{अतः} (ataḥ)}{अतः (ataḥ)}]{\hi{अतः} (ataḥ) — therefore, in formal writing}

\label{lesson:HI-C281-atah}



\begin{quote}
Before the new one: say \textquotedblleft{}so, therefore\textquotedblright{} the everyday way. (\textbf{\hi{इसलिए}}, \emph{isliye}.)
\end{quote}

\subsection*{You'll want to know: \hi{अतः}}
\textbf{\hi{अतः}} — \emph{ataḥ} — \textquotedblleft{}therefore\textquotedblright{}. It is formal. You will read it in notices, letters and the news; in speech, Hindi says \textbf{\hi{इसलिए}}.

\begin{quote}
\textbf{\hi{कल} \hi{छुट्टी} \hi{है}, \hi{अतः} \hi{दफ़्तर} \hi{बंद} \hi{रहेगा।}}
\end{quote}

\begin{quote}
\emph{kal chuṭṭī hai, ataḥ daftar band rahegā.} — \textquotedblleft{}Tomorrow is a holiday; the office will therefore stay closed.\textquotedblright{}
\end{quote}

Its last sign, the two dots \textbf{\hi{ः}}, is new, and it has the next lesson to itself. It is said as a light breath of \emph{h} after the vowel: \emph{ataḥ}.

\begin{cousinweb}
\textbf{\hi{अतः}} (\emph{ataḥ}) is a tatsama, a Sanskrit word Hindi uses unchanged. Monier-Williams gives its sense as \textquotedblleft{}from this\textquotedblright{}, and so \textquotedblleft{}hence, for this reason\textquotedblright{}: the result comes from what was said before. CDIAL has no row for it.
\end{cousinweb}

\subsection*{Your turn}
\begin{itemize}
\raggedright
  \item \emph{Say it:} \emph{ataḥ}, with the light breath at the end
  \item \emph{Read it:} \textbf{\hi{कल} \hi{छुट्टी} \hi{है}, \hi{अतः} \hi{दफ़्तर} \hi{बंद} \hi{रहेगा।}}
  \item \emph{Say it:} the same sentence the everyday way — \emph{kal chuṭṭī hai, isliye daftar band rahegā}
\end{itemize}

\begin{tcolorbox}[breakable,colback=teal!4,colframe=teal!35!black,title={Before you move on}]
What does \textbf{\hi{अतः}} mean? (\textquotedblleft{}Therefore\textquotedblright{}.) Where would you meet it? (In notices, letters and the news.) What does speech say instead? (\textbf{\hi{इसलिए}}.)
\end{tcolorbox}
\section[\texorpdfstring{\hi{ः} (ḥ)}{ः (ḥ)}]{\hi{◌ः} — two dots and a breath}

\label{lesson:HI-W281-sign-visarga}



\begin{quote}
Say \textbf{\hi{अतः}}, \textquotedblleft{}therefore\textquotedblright{}. Then find the dot of \textbf{\hi{◌ं}} in \textbf{\hi{मैं}}: one dot, above the head-line.

Today's sign has two dots, and they stand beside the line, not above it.
\end{quote}

\subsection*{Script you'll notice: \hi{◌ः}}
\textbf{\hi{◌ः}} — the \textbf{visarga}, \emph{ḥ}: a light breath of \emph{h} after the vowel before it.

Where it sits: \textbf{to the right of its syllable}, two dots one above the other. In print they hang under a short piece of head-line.

What it is made of:

\begin{itemize}
\raggedright
  \item an upper dot
  \item a lower dot
  \item a short piece of headline above them
\end{itemize}

Worked through: \textbf{\hi{त}} + \textbf{\hi{◌ः}} = \textbf{\hi{तः}}, the end of \textbf{\hi{अतः}}.

It is rare in everyday Hindi. You will meet it in formal words taken straight from Sanskrit, such as \textbf{\hi{अतः}}.

\subsection*{Writing: \hi{◌ः} — follow the numbered strip}
\hlblockfigure{\detokenize{figures/HI-W281-sign-visarga-filmstrip.pdf}}{How \hi{ः} is written, stroke by stroke}

\begin{itemize}
\raggedright
  \item \textbf{1.} start at the top of the upper dot and curve down its left side
  \item \textbf{2.} without lifting, round the bottom and up the right side to close the loop
  \item \textbf{3.} lift, then start at the top of the lower dot and curve down its left side
  \item \textbf{4.} without lifting, round the bottom and up the right side to close the loop — and only now lift
\end{itemize}

\textbf{Pen lifts: 1.} The upper dot first, then the lower.

The short piece of head-line stays grey in the strip: in a word it is part of the head-line, which you draw last, once, across the whole word.

\begin{quote}
Two strokes — four movements, one pen lift; the upper dot, then the lower, each a small loop drawn anticlockwise from its top, as \hi{ं} is; the short piece of headline the printed sign carries is left undrawn, because in a word it is part of the one headline drawn last across the whole word; cited to native writers' pen traces in HP Labs India's LipiTk Devanagari recognizer, class 60.
\end{quote}

\begin{quote}
This is one attested order, and handwriting varies by writer. Source: HP Labs India, Lipi Toolkit, Lipi Indic Character Recognizers 4.0, Devanagari recognizer, class 60 (\hi{ः}, visarga): stored online-handwriting prototypes from native writers (MIT licence, 2012).
\end{quote}

\subsection*{Your turn}
\begin{itemize}
\raggedright
  \item \emph{Look:} at these two signs, and say where each one sits
\end{itemize}

\begin{quote}
\hi{◌ं}  ·  \hi{◌ः}
\end{quote}

\begin{itemize}
\raggedright
  \item \emph{Trace it:} \hi{◌ः} three times, the upper dot first
  \item \emph{Write it:} \textbf{\hi{अतः}}, then read it aloud
\end{itemize}

\begin{tcolorbox}[breakable,colback=teal!4,colframe=teal!35!black,title={Before you move on}]
Where does \textbf{\hi{◌ः}} sit? (To the right of its syllable.) Which dot comes first? (The upper.) What does it add? (A light breath of \emph{h}.) Which word brought it? (\textbf{\hi{अतः}}.)
\end{tcolorbox}
\section[\texorpdfstring{\hi{बारिश} \hi{हो} \hi{रही} \hi{थी} … (bāriś ho rahī …)}{बारिश हो रही थी … (bāriś ho rahī …)}]{\hi{बारिश} \hi{हो} \hi{रही} \hi{थी}, \hi{इसलिए} … (bāriś ho rahī thī, isliye …) — the cause, then so}

\label{lesson:HI-C281-isliye}



\begin{quote}
Say \textquotedblleft{}because of the rain\textquotedblright{}. (\textbf{\hi{बारिश} \hi{की} \hi{वजह} \hi{से}}.) Then the two words for \textquotedblleft{}therefore\textquotedblright{}, the everyday one and the formal one. (\textbf{\hi{इसलिए}}, \textbf{\hi{अतः}}.)
\end{quote}

\subsection*{You'll want to know: cause, \hi{इसलिए}, result}
\begin{quote}
\textbf{\hi{बारिश} \hi{हो} \hi{रही} \hi{थी}, \hi{इसलिए} \hi{मैं} \hi{नहीं} \hi{गया।}}
\end{quote}

\begin{quote}
\emph{bāriś ho rahī thī, isliye maiṁ nahīṃ gayā.} — \textquotedblleft{}It was raining, so I didn't go.\textquotedblright{}
\end{quote}

\textbf{\hi{इसलिए}} stands between two whole sentences: the cause first, then the result. \textbf{\hi{क्योंकि}} turns the order round.

\noindent
\begin{tabularx}{\linewidth}{@{}>{\raggedright\arraybackslash}X>{\raggedright\arraybackslash}X@{}}
\toprule
\textbf{order} & \textbf{sentence} \\
\midrule
cause, \textbf{\hi{इसलिए}}, result & \textbf{\hi{बारिश} \hi{हो} \hi{रही} \hi{है}, \hi{इसलिए} \hi{मैं} \hi{घर} \hi{पर} \hi{हूँ।}} (\emph{bāriś ho rahī hai, isliye maiṁ ghar par hūṁ.}) \\
result, \textbf{\hi{क्योंकि}}, cause & \textbf{\hi{मैं} \hi{घर} \hi{पर} \hi{हूँ}, \hi{क्योंकि} \hi{बारिश} \hi{हो} \hi{रही} \hi{है।}} (\emph{maiṁ ghar par hūṁ, kyoṁki bāriś ho rahī hai.}) \\
\bottomrule
\end{tabularx}

In formal writing, \textbf{\hi{अतः}} takes the place of \textbf{\hi{इसलिए}}.

\subsection*{You'll want to know: \hi{तो}, so then}
\begin{quote}
\textbf{\hi{बारिश} \hi{हो} \hi{रही} \hi{है}? \hi{तो} \hi{छाता} \hi{लीजिए।}}
\end{quote}

\begin{quote}
\emph{bāriś ho rahī hai? to chātā lījie.} — \textquotedblleft{}It's raining? Then take an umbrella.\textquotedblright{}
\end{quote}

\textbf{\hi{तो}} (\emph{to}), \textquotedblleft{}so, then\textquotedblright{}, takes up what was said and draws the next step from it. It opens a reply:

\begin{itemize}
\raggedright
  \item \textbf{\hi{तो} \hi{चलिए।}} (\emph{to calie.}) — So let's go.
  \item \textbf{\hi{ठीक} \hi{है}, \hi{तो} \hi{कल} \hi{मिलते} \hi{हैं।}} (\emph{ṭhīk hai, to kal milte haiṁ.}) — All right, so see you tomorrow.
\end{itemize}

\textbf{\hi{तो}} has a bigger job still, after \textbf{\hi{अगर}}, \textquotedblleft{}if\textquotedblright{}; it comes when the book reaches conditions.

\subsection*{Your turn}
Say these aloud:

\begin{itemize}
\raggedright
  \item \textquotedblleft{}It's raining, so I'm at home\textquotedblright{} — \emph{bāriś ho rahī hai, isliye maiṁ ghar par hūṁ}
  \item \textquotedblleft{}All right, so see you tomorrow\textquotedblright{} — \emph{ṭhīk hai, to kal milte haiṁ}
  \item \textquotedblleft{}You're ill? Then you should rest\textquotedblright{} — \emph{tum bīmār ho? to tumhẽ ārām karnā cāhiye}
\end{itemize}

\begin{tcolorbox}[breakable,colback=teal!4,colframe=teal!35!black,title={Before you move on}]
Where does \textbf{\hi{इसलिए}} go? (Between the cause and the result.) And \textbf{\hi{क्योंकि}}? (After the result, before the cause.) \textquotedblleft{}So let's go\textquotedblright{}? (\textbf{\hi{तो} \hi{चलिए}}.) What does formal writing put for \textbf{\hi{इसलिए}}? (\textbf{\hi{अतः}}.)
\end{tcolorbox}
\section[(review)]{Review — advice, and so}

\label{lesson:HI-R281-advice-and-reasons}



\begin{quote}
Say \textquotedblleft{}you should rest\textquotedblright{}. (\emph{tumhẽ ārām karnā cāhiye}.) Then read \textbf{\hi{अतः}} aloud, and write its last sign: the upper dot, then the lower.
\end{quote}

\subsection*{Your turn}
\begin{itemize}
\raggedright
  \item \emph{Say it:} You should drink water. I had to go to the office.
  \item \emph{Say it:} It was raining, so we stayed at home.
  \item \emph{Say it:} You're ill? Then you should rest.
  \item \emph{Say it:} I think we will have to leave early.
  \item \emph{Read it:} \textbf{\hi{कल} \hi{छुट्टी} \hi{है}, \hi{अतः} \hi{दफ़्तर} \hi{बंद} \hi{रहेगा।}} Then say it the everyday way.
\end{itemize}

\begin{tcolorbox}[breakable,colback=teal!4,colframe=teal!35!black,title={Before you move on}]
\textquotedblleft{}Should\textquotedblright{}? (\textbf{\hi{चाहिए}}, after a verb in \textbf{-\hi{ना}}.) \textquotedblleft{}Had to\textquotedblright{}? (\textbf{\hi{पड़ा}}.) \textquotedblleft{}So\textquotedblright{}, between two sentences? (\textbf{\hi{इसलिए}}, or in writing \textbf{\hi{अतः}}.) \textquotedblleft{}So then\textquotedblright{}, opening a reply? (\textbf{\hi{तो}}.) Where does \textbf{\hi{◌ः}} sit? (To the right of its syllable.)
\end{tcolorbox}
